% !TEX root = paper-nifty.tex

% In this paper, we presented net2vec, a large-scale feature learning framework for networks motivated by recent advances in representation learning. Our framework, composed of two independent algorithms for generating feature representations for nodes and edges respectively. 

% The first component, node2vec, is based on random walks over the underlying network structure. We see that biasing our walks to account for the distributed influence of nodes as well as tuning our search procedure to move inward (breadth-first) or outward (depth-first) is suited for different kinds of similarity inference (homophilic/structural equivalence). The level to which each type of similarity inference is of interest is governed by two tunable hyper-parameters, that can either by learned in a semi-supervised fashion or set based on domain expertize. In our second component, edge2vec, we generated latent representations for pairs of nodes based on binary operations over node representations. Empirically, we observed significant gains  over benchmark algorithms on the multi-label classification tasks for nodes and the link prediction and signed edge prediction tasks for edges on several real-world datasets.

% In the future, we would like to work towards interpreting and correlating the features optimized using single-hidden layer neural-networks to heuristic measures commonly used in network science to get a better understanding. Extending these methods to more specialized kinds of networks such as heterogenous information networks, bipartitite graphs or networks with features for nodes/edges is a potential direction to extend this work. With respect to the latter, a recent work \cite{babis-www15} has explored the inverse problem of inferring binary latent factors through MCMC methods under an assumption that networks are homophilous. Finally, it would be interesting to see the potential implications of intermediate representational learning methods such as node2vec for fast training of end-to-end deep learning methods.  

In this paper, we presented \nodevec, a scalable algorithmic framework for feature learning in networks. Our algorithm simulates biased random walks over the underlying network. The bias provides flexibility in exploring neighborhoods. This flexibility in search is important in discovering a mixture of equivalences corresponding to homophily and structural equivalence. In addition to good predicitive accuracy, \nodevec is also computationally efficient and can scale to large real-world networks.

There are several directions for future work. Our algorithm assumes a general setting of an (un)directed, (un)weighted network. We hope to extend it to specialized networks such as heterogeneous information networks, networks with explicit domain features for nodes and edges and signed-edge networks which provide additional information about the domain. A large part of our analysis concerning node equivalences in networks is general, and it would be interesting to extend these ideas to design network-aware supervised feature learning methods such as end-to-end deep learning architectures.
